% AUTHORED
\section{Limits of linearity (XOR)}\label{sec:02}
\conceptbadge{02}

The perceptron draws one straight decision boundary. That works only when the
classes are linearly separable. XOR exposes the limit. For two binary inputs,
XOR returns $1$ when exactly one input is $1$:
\[
\begin{array}{ccc}
 x_1 & x_2 & y\\
 0 & 0 & 0\\
 0 & 1 & 1\\
 1 & 0 & 1\\
 1 & 1 & 0
\end{array}
\]
The positive examples sit on opposite corners of the unit square. No single
line can put both positive corners on one side while keeping both negative
corners on the other side.

\subsection*{More linear layers are not enough}
Suppose you stack two affine maps but use no nonlinearity:
\[
  h = xW_1 + b_1,
  \qquad
  y = hW_2 + b_2.
\]
Substitute the first equation into the second:
\[
  y = x(W_1W_2) + (b_1W_2 + b_2).
\]
The composition is still just another affine map. No matter how many linear
layers you stack, without a nonlinearity the whole network collapses to one
linear layer.

\subsection*{Hidden features}
A nonlinear hidden layer can divide the input space into regions, and a
final linear unit can then combine those regions. For XOR, one hidden unit can detect
whether at least one input is active, while another can detect whether both are
active. A final unit can combine those features as
\[
  \text{at-least-one} - \text{both}.
\]
This is the core idea of representation learning: the early layers do not merely
classify; they transform the data into features where classification is simpler.

\subsection*{Why this matters later}
The transformer's feed-forward network is an MLP applied at every token
position. Attention mixes information across positions, but the MLP supplies
per-token nonlinear feature transformation. The XOR example is tiny, yet it
captures the reason that nonlinearity is not optional.

\subsection*{Exercise connection}
The exercise gives you a small two-layer network skeleton for XOR. You are not
training it yet; the point is to see how a hidden representation can make a
problem linearly separable for the output unit.
